% =====================================================================
\chapter{Estimación puntual e intervalos de confianza}
\label{tema:4}
% =====================================================================

El \cref{tema:2} enseñó a juzgar un estimador; este tema enseña a construirlo. Se presentan los dos métodos generales de estimación puntual, el método de los momentos y el de máxima verosimilitud, y se aplican a los modelos de Bernoulli, Poisson, normal y exponencial. Después se pasa de la estimación puntual a la estimación por intervalos: intervalos de confianza para la media con varianza conocida y desconocida, para proporciones y cálculo del tamaño muestral. Se sigue el capítulo~9 de Wasserman~\cite{wasserman}, «Parametric Inference» (método de los momentos y máxima verosimilitud), las clases~12, 13 y~18 de Stats~200~\cite{stats200}, las lecturas~10b, 22 y~23a del curso 18.05 del MIT~\cite{mit1805} y el capítulo~8 de OpenStax~\cite{openstax}.

\section{El método de los momentos}

Sea $X_1,\dots,X_n$ una muestra \iid\ de un modelo $F_\theta$ con $\theta=(\theta_1,\dots,\theta_k)$. El \emph{momento poblacional} de orden $j$ es $m_j(\theta)=\E_\theta[X^{j}]$ y el \emph{momento muestral} de orden $j$ es $\widehat m_j=\frac1n\sum_iX_i^{j}$. Por la ley de los grandes números, $\widehat m_j\convp m_j(\theta)$.

\begin{definicion}[Estimador por el método de los momentos]\label{def:t4-momentos}
El \emph{estimador por el método de los momentos} $\widehat\theta$ es la solución del sistema de $k$ ecuaciones
\[
m_j(\widehat\theta)=\widehat m_j,\qquad j=1,\dots,k,
\]
es decir, el valor del parámetro con el que los primeros $k$ momentos del modelo coinciden con los de la muestra~\cite[cap.~9]{wasserman}.
\end{definicion}

\begin{ejemplo}[Momentos en los modelos básicos]\label{ej:t4-momentos}
\leavevmode
\begin{enumerate}
  \item Bernoulli: $\E X=p$, luego $\widehat p=\Xbar$.
  \item Poisson: $\E X=\lambda$, luego $\widehat\lambda=\Xbar$.
  \item Exponencial de tasa $\lambda$: $\E X=1/\lambda$, luego $\widehat\lambda=1/\Xbar$.
  \item Normal: $\E X=\mu$ y $\E X^{2}=\sigma^{2}+\mu^{2}$. De $\widehat\mu=\Xbar$ y $\widehat\sigma^{2}+\widehat\mu^{2}=\frac1n\sum_iX_i^{2}$ se obtiene $\widehat\sigma^{2}=\frac1n\sum_iX_i^{2}-\Xbar^{2}=\frac1n\sum_i(X_i-\Xbar)^{2}$, la varianza con denominador $n$.
  \item Uniforme en $[0,\theta]$: $\E X=\theta/2$, luego $\widehat\theta=2\Xbar$ (\cref{eje:t2-uniforme}).
\end{enumerate}
\end{ejemplo}

El método es sencillo y produce estimadores consistentes en condiciones generales, pero no siempre los mejores y, a veces, incompatibles con los datos (\cref{eje:t4-uniforme-ab}).

\section{Máxima verosimilitud}

\begin{definicion}[Verosimilitud y estimador de máxima verosimilitud]\label{def:t4-emv}
Sean $x_1,\dots,x_n$ los datos observados y $f(x\mid\theta)$ la función de probabilidad (caso discreto) o la densidad (caso continuo) del modelo. La \emph{función de verosimilitud} y la \emph{log-verosimilitud} son
\[
\Lik(\theta)=\prod_{i=1}^{n}f(x_i\mid\theta),\qquad \ell(\theta)=\ln\Lik(\theta)=\sum_{i=1}^{n}\ln f(x_i\mid\theta).
\]
El \emph{estimador de máxima verosimilitud} (EMV) es el valor $\widehat\theta=\argmax_{\theta\in\Theta}\Lik(\theta)$~\cite[cap.~9]{wasserman}.
\end{definicion}

La verosimilitud es la probabilidad (o densidad) de los datos observados vista como función del parámetro, con los datos fijos. Como el logaritmo es estrictamente creciente, maximizar $\Lik$ equivale a maximizar $\ell$, que es una suma y se deriva con más facilidad. El procedimiento habitual es resolver la \emph{ecuación de verosimilitud} $\ell'(\theta)=0$ y comprobar que la solución es un máximo global.

\begin{teorema}[EMV de Bernoulli]\label{teo:t4-emv-bernoulli}
Si $x_1,\dots,x_n\in\{0,1\}$ son observaciones \iid\ de $\Ber(p)$ y $k=\sum_ix_i$ es el número de unos, con $0<k<n$, el EMV es $\widehat p=k/n=\xbar$.
\end{teorema}

\begin{proof}
La función de probabilidad es $f(x\mid p)=p^{x}(1-p)^{1-x}$ para $x\in\{0,1\}$. Entonces $\Lik(p)=p^{k}(1-p)^{n-k}$ y, para $0<p<1$,
\[
\ell(p)=k\ln p+(n-k)\ln(1-p),\qquad \ell'(p)=\frac kp-\frac{n-k}{1-p}.
\]
La ecuación $\ell'(p)=0$ equivale a $k(1-p)=(n-k)p$, es decir, $p=k/n$. Como $\ell''(p)=-k/p^{2}-(n-k)/(1-p)^{2}<0$, $\ell$ es estrictamente cóncava en $(0,1)$ y el punto crítico es el máximo global. (En los extremos, $\Lik(0)=\Lik(1)=0$ cuando $0<k<n$.)
\end{proof}

Si $k=0$ o $k=n$, el máximo está en $\widehat p=0$ o $\widehat p=1$, respectivamente, y la fórmula $\widehat p=k/n$ sigue siendo válida.

\begin{teorema}[EMV de la normal]\label{teo:t4-emv-normal}
Si $x_1,\dots,x_n$ son observaciones \iid\ de $\Nor(\mu,\sigma^{2})$, no todas iguales, el EMV de $(\mu,\sigma^{2})$ es
\[
\widehat\mu=\xbar,\qquad\widehat\sigma^{2}=\frac1n\sum_{i=1}^{n}(x_i-\xbar)^{2}.
\]
\end{teorema}

\begin{proof}
Con $v=\sigma^{2}>0$, la log-verosimilitud es
\[
\ell(\mu,v)=-\frac n2\ln(2\pi)-\frac n2\ln v-\frac{1}{2v}\sum_{i=1}^{n}(x_i-\mu)^{2}.
\]
\emph{Paso 1 ($\mu$).} Para cada $v$ fijo, maximizar $\ell$ en $\mu$ equivale a minimizar $\sum_i(x_i-\mu)^{2}$, cuyo mínimo está en $\mu=\xbar$ (\cref{prop:t1-extremal}). Por tanto $\ell(\mu,v)\le\ell(\xbar,v)$ para todo $(\mu,v)$.

\emph{Paso 2 ($v$).} Sea $Q=\sum_i(x_i-\xbar)^{2}>0$ y $h(v)=\ell(\xbar,v)=c-\frac n2\ln v-\frac{Q}{2v}$. Entonces $h'(v)=-\frac{n}{2v}+\frac{Q}{2v^{2}}=\frac{Q-nv}{2v^{2}}$, positiva para $v<Q/n$ y negativa para $v>Q/n$. Luego $h$ alcanza su máximo en $v=Q/n$.

\emph{Paso 3.} Combinando: $\ell(\mu,v)\le\ell(\xbar,v)\le\ell(\xbar,Q/n)$ para todo $(\mu,v)$.
\end{proof}

El EMV de $\sigma^{2}$ coincide con el estimador por momentos y es sesgado (\cref{cor:t2-sesgo}); por eso, para datos, se usa casi siempre $S^{2}$.

\begin{ejemplo}[Un caso sin derivadas]\label{ej:t4-uniforme}
Si $x_1,\dots,x_n$ son observaciones de $\Unif(0,\theta)$, entonces $f(x\mid\theta)=1/\theta$ si $0\le x\le\theta$ y $0$ en otro caso, luego
\[
\Lik(\theta)=\begin{cases}\theta^{-n}&\text{si }\theta\ge\max_ix_i,\\0&\text{si }\theta<\max_ix_i.\end{cases}
\]
$\Lik$ es decreciente en $[\max_ix_i,\infty)$, de modo que $\widehat\theta=\max_ix_i$. La ecuación $\ell'(\theta)=0$ no tiene solución: el máximo está en un extremo del dominio.
\end{ejemplo}

\begin{observacion}[Propiedades del EMV]\label{obs:t4-propiedades}
Se enuncian sin demostración, como ampliación~\cite[cap.~9]{wasserman}: (a) \emph{invariancia}: si $\widehat\theta$ es el EMV de $\theta$, el EMV de $g(\theta)$ es $g(\widehat\theta)$; por ejemplo, el EMV de $\sigma$ es $\sqrt{\widehat\sigma^{2}}$; (b) en condiciones de regularidad (que no se cumplen en el \cref{ej:t4-uniforme}), el EMV es consistente y aproximadamente normal para $n$ grande, con un error estándar que se puede estimar a partir de la curvatura de $\ell$ en el máximo.
\end{observacion}

\section{Intervalos de confianza}

Una estimación puntual no informa de su precisión. Un intervalo de confianza acompaña la estimación de un margen de error con una garantía probabilística.

\begin{definicion}[Intervalo de confianza]\label{def:t4-ic}
Sea $0<\alpha<1$. Un \emph{intervalo de confianza} de nivel $1-\alpha$ para $\theta$ es un intervalo $[L,U]$ cuyos extremos $L=L(X_1,\dots,X_n)$ y $U=U(X_1,\dots,X_n)$ son estadísticos y tal que
\[
\PP_\theta(L\le\theta\le U)\ge1-\alpha\qquad\text{para todo }\theta\in\Theta.
\]
La probabilidad $\PP_\theta(L\le\theta\le U)$ se llama \emph{cobertura}~\cite[cap.~6]{wasserman}.
\end{definicion}

\begin{observacion}[Interpretación]\label{obs:t4-interpretacion}
Lo aleatorio es el intervalo, no el parámetro. Si se repitiera muchas veces el experimento y se calculara cada vez el intervalo de nivel $95\,\%$, aproximadamente el $95\,\%$ de esos intervalos contendría el verdadero $\theta$. Una vez observados los datos, el intervalo concreto (por ejemplo, $[0{,}79;\allowbreak\,0{,}84]$) contiene o no contiene a $\theta$; no es correcto decir que $\theta$ está en él «con probabilidad $0{,}95$» en el sentido frecuentista~\cite[lectura~23a]{mit1805}.
\end{observacion}

El método general para construir intervalos usa un \emph{pivote}: una función $Q(X_1,\dots,X_n;\theta)$ cuya distribución no depende de $\theta$.

\begin{teorema}[Intervalo para la media con $\sigma$ conocida]\label{teo:t4-ic-z}
Sean $X_1,\dots,X_n\sim\Nor(\mu,\sigma^{2})$ \iid\ con $\sigma$ conocida. El intervalo
\[
\Bigl[\Xbar-z_{1-\alpha/2}\frac{\sigma}{\sqrt n},\ \Xbar+z_{1-\alpha/2}\frac{\sigma}{\sqrt n}\Bigr]
\]
tiene cobertura exactamente $1-\alpha$. Si los datos no son normales pero $n$ es grande, la cobertura es aproximadamente $1-\alpha$.
\end{teorema}

\begin{proof}
El pivote $Z=(\Xbar-\mu)/(\sigma/\sqrt n)$ tiene distribución $\Nor(0,1)$ para todo $\mu$ (\cref{teo:t3-suma-normales,prop:t3-tipificacion}). Con $c=z_{1-\alpha/2}$, $\PP(-c\le Z\le c)=\Phi(c)-\Phi(-c)=2\Phi(c)-1=1-\alpha$. Despejando $\mu$ en las desigualdades,
\[
-c\le\frac{\Xbar-\mu}{\sigma/\sqrt n}\le c\iff\Xbar-c\frac{\sigma}{\sqrt n}\le\mu\le\Xbar+c\frac{\sigma}{\sqrt n}.
\]
Los dos sucesos son el mismo y tienen la misma probabilidad. Para datos no normales se aplica el mismo argumento con la aproximación del \cref{teo:t3-tcl}.
\end{proof}

La cantidad $E=z_{1-\alpha/2}\,\sigma/\sqrt n$ es el \emph{margen de error} (la semiamplitud del intervalo). Los valores más usados son $z_{0{,}95}=1{,}645$ (nivel $90\,\%$), $z_{0{,}975}=1{,}960$ ($95\,\%$) y $z_{0{,}995}=2{,}576$ ($99\,\%$).

\begin{corolario}[Tamaño muestral]\label{cor:t4-tamano}
Para que el margen de error del \cref{teo:t4-ic-z} sea a lo sumo $E$ basta tomar
\[
n\ge\Bigl(\frac{z_{1-\alpha/2}\,\sigma}{E}\Bigr)^{2},
\]
redondeando al entero superior.
\end{corolario}

\begin{proof}
$z_{1-\alpha/2}\sigma/\sqrt n\le E\iff\sqrt n\ge z_{1-\alpha/2}\sigma/E$, y se elevan al cuadrado ambos miembros, que son positivos.
\end{proof}

\section{Intervalo para la media con \texorpdfstring{$\sigma$}{sigma} desconocida}

En la práctica $\sigma$ casi nunca se conoce. Al sustituirla por $S$ hay que reemplazar la normal por la $t$ de Student.

\begin{teorema}[Intervalo $t$]\label{teo:t4-ic-t}
Sean $X_1,\dots,X_n\sim\Nor(\mu,\sigma^{2})$ \iid\ con $\sigma$ desconocida y $n\ge2$. El intervalo
\[
\Bigl[\Xbar-t_{n-1;1-\alpha/2}\frac{S}{\sqrt n},\ \Xbar+t_{n-1;1-\alpha/2}\frac{S}{\sqrt n}\Bigr]
\]
tiene cobertura exactamente $1-\alpha$.
\end{teorema}

\begin{proof}
Por el \cref{cor:t3-t}, el pivote $T=(\Xbar-\mu)/(S/\sqrt n)$ tiene distribución $t_{n-1}$, que no depende de $\mu$ ni de $\sigma$. Con $c=t_{n-1;1-\alpha/2}$, la simetría de la $t$ (\cref{eje:t3-t-simetrica}) da $\PP(-c\le T\le c)=1-\alpha$, y se despeja $\mu$ como en el \cref{teo:t4-ic-z}.
\end{proof}

Como $t_{n-1;1-\alpha/2}>z_{1-\alpha/2}$, el intervalo $t$ es más ancho: refleja la incertidumbre adicional por estimar $\sigma$. Para $n$ grande (digamos $n\ge30$ o más) ambos cuantiles casi coinciden y, por el teorema central del límite, el intervalo es aproximadamente válido aunque los datos no sean normales.

\begin{ejemplo}\label{ej:t4-t}
Se miden los tiempos de respuesta (en segundos) de un sistema en $n=10$ consultas y se obtiene $\xbar=2{,}40$ y $s=0{,}50$. Con $t_{9;0{,}975}=2{,}262$, el intervalo de nivel $95\,\%$ es $2{,}40\pm2{,}262\cdot0{,}50/\sqrt{10}=2{,}40\pm0{,}358$, es decir, $[2{,}04;\allowbreak\ 2{,}76]$. Si se hubiera usado $1{,}96$, el margen habría sido $0{,}310$, un $13\,\%$ menor.
\end{ejemplo}

\section{Intervalos para proporciones}

Si $X_1,\dots,X_n\sim\Ber(p)$, el \cref{cor:t3-binomial} da $\widehat p\approx\Nor\bigl(p,p(1-p)/n\bigr)$. El error estándar depende de $p$, desconocido, y se estima por $\sqrt{\widehat p(1-\widehat p)/n}$.

\begin{definicion}[Intervalo de Wald para una proporción]\label{def:t4-wald}
El intervalo de confianza aproximado de nivel $1-\alpha$ para $p$ es
\[
\widehat p\pm z_{1-\alpha/2}\sqrt{\frac{\widehat p(1-\widehat p)}{n}}.
\]
Como $p(1-p)\le1/4$, el intervalo \emph{conservador} $\widehat p\pm z_{1-\alpha/2}/(2\sqrt n)$ es siempre al menos tan ancho~\cite[§8.3]{openstax}.
\end{definicion}

La aproximación es buena si $n\widehat p$ y $n(1-\widehat p)$ son ambos al menos del orden de $10$; con muestras pequeñas o con $\widehat p$ cercano a $0$ o a $1$ la cobertura real es inferior a la nominal. Existen correcciones sencillas, como el método de «más cuatro» de~\cite[§8.3]{openstax} (sumar dos éxitos y dos fracasos antes de calcular) o el intervalo de Wilson, que quedan fuera de este módulo.

\begin{corolario}[Tamaño muestral para una proporción]\label{cor:t4-tamano-p}
Para que el margen de error sea a lo sumo $E$ sea cual sea $p$, basta $n\ge z_{1-\alpha/2}^{2}/(4E^{2})$. Si se dispone de una estimación previa $p^{*}$, basta $n\ge z_{1-\alpha/2}^{2}\,p^{*}(1-p^{*})/E^{2}$.
\end{corolario}

\begin{proof}
Es el \cref{cor:t4-tamano} con $\sigma=\sqrt{p(1-p)}$, acotado por $1/2$ o sustituido por $\sqrt{p^{*}(1-p^{*})}$.
\end{proof}

\begin{ejemplo}[Exactitud en un benchmark]\label{ej:t4-benchmark}
Un modelo acierta $820$ de $1000$ preguntas. Entonces $\widehat p=0{,}82$, $\widehat{\ee}=\sqrt{0{,}82\cdot0{,}18/1000}\approx0{,}01215$ y el intervalo de nivel $95\,\%$ es $0{,}82\pm1{,}96\cdot0{,}01215=0{,}82\pm0{,}0238$, es decir, $[0{,}796;\allowbreak\ 0{,}844]$. Para reducir el margen a un punto porcentual ($E=0{,}01$) harían falta $1{,}96^{2}\cdot0{,}82\cdot0{,}18/0{,}01^{2}\approx5671$ preguntas.
\end{ejemplo}

\begin{observacion}[Errores frecuentes]\label{obs:t4-errores}
(a) Interpretar un intervalo de confianza del $95\,\%$ como «hay un $95\,\%$ de probabilidad de que $\theta$ esté en este intervalo concreto» o como «el $95\,\%$ de los datos está en el intervalo» (\cref{obs:t4-interpretacion}).
(b) Usar $z$ en lugar de $t$ con $n$ pequeño y $\sigma$ estimada.
(c) Olvidar que el margen de error solo recoge el error de muestreo: no protege contra sesgos (por ejemplo, preguntas del benchmark que estaban en los datos de entrenamiento).
(d) Para comparar dos modelos, mirar solo si sus intervalos se solapan; dos intervalos que se solapan ligeramente pueden corresponder a una diferencia significativa (\cref{tema:5}).
(e) En el EMV, derivar sin comprobar que el punto crítico es un máximo, o derivar cuando el soporte depende del parámetro (\cref{ej:t4-uniforme}).
\end{observacion}

\begin{observacion}[Conexión con el curso]
El módulo~1 obtiene el EMV de un modelo de bigramas, que generaliza el \cref{teo:t4-emv-bernoulli} a distribuciones categóricas: la frecuencia relativa. Entrenar un modelo de lenguaje minimizando la entropía cruzada es exactamente maximizar la log-verosimilitud del corpus (módulos~1, 5 y~8). En los módulos~22 y~23 se construyen intervalos de confianza para la exactitud en un benchmark (\cref{ej:t4-benchmark}) y la regla del~3 (\cref{eje:t4-regla3}) para acotar la tasa de fallos graves cuando no se ha observado ninguno; en el módulo~16, los exponentes de las leyes de escala se estiman con su incertidumbre.
\end{observacion}

\begin{lecturas}
\item Wasserman~\cite{wasserman}, cap.~9 «Parametric Inference» (método de los momentos y máxima verosimilitud) y cap.~6 (intervalos de confianza).
\item Stats~200~\cite{stats200}, clases~12 «Parametric models, method of moments estimation», 13 «Maximum likelihood estimation» y 18 «Confidence intervals».
\item Orloff y Bloom~\cite{mit1805}, lecturas~10b «Maximum Likelihood Estimates», 22 «Confidence intervals based on normal data», 23a «Confidence Intervals: Three Views» y 23b «Confidence Intervals for the Mean of Non-normal Data».
\item OpenStax~\cite{openstax}, cap.~8: §8.1 «A Single Population Mean using the Normal Distribution», §8.2 «A Single Population Mean using the Student t Distribution» y §8.3 «A Population Proportion».
\item Rice~\cite{rice}, cap.~8 «Estimation of Parameters and Fitting of Probability Distributions».
\end{lecturas}

\section*{Ejercicios del tema 4}
\addcontentsline{toc}{section}{Ejercicios del tema 4}

\begin{ejercicio}{1}{eje:t4-interpretacion}
Un estudio concluye que un nuevo método de ajuste fino mejora la exactitud media en $15$ puntos, con un intervalo de confianza del $95\,\%$ de $[10,20]$ puntos. Indique cuáles de estas afirmaciones se deducen del intervalo: (i) hay una probabilidad del $95\,\%$ de que la mejora real esté entre $10$ y $20$; (ii) la probabilidad de que el método mejore la exactitud es al menos del $95\,\%$; (iii) el intervalo da una estimación de la mejora real; (iv) la hipótesis de que no hay mejora es probablemente falsa; (v) si se repitiera el estudio muchas veces, alrededor del $95\,\%$ de los intervalos contendría la mejora real; (vi) un contraste bilateral al $5\,\%$ rechazaría la hipótesis de mejora nula.
\fuente{adaptado de Orloff y Bloom, MIT 18.05~\cite{mit1805}, \emph{problem set}~10, problema~5(c).}
\end{ejercicio}

\begin{ejercicio}{2}{eje:t4-momentos-basico}
(a) Los tiempos entre peticiones a un servidor, en segundos, son $0{,}5;\allowbreak\ 1{,}2;\allowbreak\ 0{,}3;\allowbreak\ 2{,}0$. Suponiendo un modelo exponencial, estime $\lambda$ por el método de los momentos. (b) Para los datos $2,4,4,5,7$ y un modelo normal, obtenga las estimaciones de $\mu$ y $\sigma^{2}$ por el método de los momentos y compárelas con $s^{2}$.
\end{ejercicio}

\begin{ejercicio}{2}{eje:t4-geometrica}
Sean $X_1,\dots,X_n$ \iid\ con distribución geométrica en $\{1,2,3,\dots\}$, $\PP(X=x)=p(1-p)^{x-1}$, cuya media es $1/p$. (a) Obtenga el estimador de $p$ por el método de los momentos. (b) Obtenga el EMV. (c) Calcule ambos con los datos $2,1,4,3,1,1$ (número de intentos hasta que un agente resuelve una tarea).
\fuente{adaptado de Stats~200, Stanford~\cite{stats200}, \emph{homework}~5, problema~1.}
\end{ejercicio}

\begin{ejercicio}{2}{eje:t4-gamma}
Una variable positiva tiene densidad $f(x\mid a)=\dfrac{a^{5}}{4!}\,x^{4}e^{-ax}$ para $x>0$, con $a>0$. (a) Obtenga el EMV de $a$ a partir de $n$ observaciones \iid. (b) Calcúlelo para los datos $2{,}0;\allowbreak\ 3{,}0;\allowbreak\ 2{,}5;\allowbreak\ 4{,}5$.
\fuente{adaptado de Orloff y Bloom, MIT 18.05~\cite{mit1805}, \emph{problem set}~6, problema~1.}
\end{ejercicio}

\begin{ejercicio}{2}{eje:t4-uniforme-ab}
Los datos $1{,}2;\allowbreak\ 2{,}1;\allowbreak\ 1{,}3;\allowbreak\ 10{,}5;\allowbreak\ 5$ proceden de una distribución $\Unif(a,b)$. (a) Obtenga el EMV de $a$ y $b$ razonando directamente sobre la verosimilitud, sin derivar. (b) Obtenga los estimadores por el método de los momentos, sabiendo que la media es $(a+b)/2$ y la varianza $(b-a)^{2}/12$. (c) ¿Son compatibles con los datos los estimadores de (b)?
\fuente{adaptado de Orloff y Bloom, MIT 18.05~\cite{mit1805}, \emph{problem set}~6, problema~3.}
\end{ejercicio}

\begin{ejercicio}{2}{eje:t4-poisson-datos}
El número de alucinaciones detectadas en ocho documentos generados es $2,0,3,1,4,2,1,3$. Suponiendo un modelo de Poisson, (a) estime $\lambda$; (b) construya un intervalo aproximado del $95\,\%$ usando que $\Var(\Xbar)=\lambda/n$.
\end{ejercicio}

\begin{ejercicio}{2}{eje:t4-soda}
Una máquina llena latas con volúmenes $\Nor(\mu,\sigma^{2})$. Se miden nueve latas (en mL): $352,351,361,353,352,358,360,358,359$. (a) Con $\sigma=3$ conocida, construya los intervalos de confianza del $95\,\%$ y del $98\,\%$ para $\mu$. (b) Repita el apartado suponiendo $\sigma$ desconocida y compare.
\fuente{adaptado de Orloff y Bloom, MIT 18.05~\cite{mit1805}, \emph{problem set}~10, problema~2.}
\end{ejercicio}

\begin{ejercicio}{2}{eje:t4-alturas}
Se mide la estatura de $20$ personas, que se suponen normales, y se obtiene $\xbar=69{,}55$ y $s^{2}=14{,}26$ (en pulgadas y pulgadas cuadradas). (a) Construya un intervalo $t$ del $90\,\%$ para $\mu$. (b) Suponiendo conocida $\sigma=3{,}77$, construya el intervalo $z$ del $90\,\%$. (c) En la situación de (b), ¿cuántas personas hay que medir para que el intervalo del $90\,\%$ tenga una amplitud total de una pulgada?
\fuente{adaptado de Orloff y Bloom, MIT 18.05~\cite{mit1805}, \emph{problem set}~11, problema~1.}
\end{ejercicio}

\begin{ejercicio}{2}{eje:t4-openstax-z}
(a) Una población tiene $\sigma=3$; una muestra de $n=36$ da $\xbar=68$. Construya el intervalo del $90\,\%$ para $\mu$. (b) Si $\sigma=15$, ¿qué tamaño muestral hace falta para que el margen de error del $95\,\%$ sea $2$?
\fuente{adaptado de OpenStax, \emph{Introductory Statistics}~2e~\cite{openstax}, §8.1, ejemplos~8.2 y~8.7.}
\end{ejercicio}

\begin{ejercicio}{2}{eje:t4-proporcion}
(a) En una encuesta a $500$ personas, $421$ declaran usar un asistente de IA a diario. Construya el intervalo del $95\,\%$ para la proporción. (b) ¿Qué tamaño muestral garantiza un margen de error del $3\,\%$ con confianza del $90\,\%$, sin información previa sobre $p$?
\fuente{adaptado de OpenStax, \emph{Introductory Statistics}~2e~\cite{openstax}, §8.3, ejemplos~8.10 y~8.14.}
\end{ejercicio}

\begin{ejercicio}{2}{eje:t4-margen}
Se usa el intervalo conservador $\widehat p\pm z_{1-\alpha/2}/(2\sqrt n)$. (a) ¿Cuántas preguntas debe tener un benchmark para que el margen sea $0{,}01$ al $95\,\%$? Hágalo con $z=1{,}96$ y con la aproximación $z\approx2$. (b) ¿Y al $80\,\%$? (c) Una encuesta anuncia un margen de error de $\pm4$ puntos al $95\,\%$; estime cuántas personas se encuestaron.
\fuente{adaptado de Orloff y Bloom, MIT 18.05~\cite{mit1805}, problemas de la clase~23, pregunta conceptual~1 y problema~1.}
\end{ejercicio}

\begin{ejercicio}{3}{eje:t4-emv-poisson}
Sean $x_1,\dots,x_n$ observaciones \iid\ de $\Pois(\lambda)$ con $\sum_ix_i>0$. Demuestre que el EMV de $\lambda$ es $\xbar$, comprobando que es un máximo global, y que coincide con el estimador por momentos.
\end{ejercicio}

\begin{ejercicio}{3}{eje:t4-emv-expo}
Sean $x_1,\dots,x_n>0$ observaciones \iid\ de $\Expo(\lambda)$, con densidad $\lambda e^{-\lambda x}$. (a) Demuestre que el EMV de $\lambda$ es $1/\xbar$. (b) Reparametrice el modelo por la media $\beta=1/\lambda$ y demuestre directamente que el EMV de $\beta$ es $\xbar$, de acuerdo con la propiedad de invariancia.
\end{ejercicio}

\begin{ejercicio}{3}{eje:t4-laplace}
La distribución de Laplace con parámetro de posición $\mu$ tiene densidad $f(x\mid\mu)=\frac12e^{-\abs{x-\mu}}$. Demuestre que el EMV de $\mu$ es la mediana muestral. (\emph{Indicación:} \cref{eje:t1-mediana-min}.) Compare con el modelo normal, cuyo EMV de la media es $\xbar$.
\end{ejercicio}

\begin{ejercicio}{3}{eje:t4-regla3}
Se prueban $n$ entradas adversarias independientes contra un modelo y ninguna provoca un fallo grave. Sea $p$ la probabilidad de fallo en cada entrada. (a) Demuestre que el conjunto de valores de $p$ para los que observar cero fallos tiene probabilidad al menos $\alpha$ es $[0,p_{\max}]$ con $p_{\max}=1-\alpha^{1/n}$. (b) Demuestre que, para $\alpha=0{,}05$, $p_{\max}\le-\ln(0{,}05)/n\approx3/n$ («regla del~3»). (c) Calcule $p_{\max}$ y $3/n$ para $n=300$.
\end{ejercicio}

\begin{ejercicio}{3}{eje:t4-conservador}
Sea $\widehat p$ la proporción muestral de $n$ ensayos de Bernoulli. Usando la aproximación normal $\widehat p\approx\Nor(p,p(1-p)/n)$, demuestre que el intervalo conservador $\widehat p\pm z_{1-\alpha/2}/(2\sqrt n)$ tiene cobertura aproximada al menos $1-\alpha$ para todo $p\in(0,1)$, con igualdad solo si $p=1/2$. ¿Qué ocurre con la cobertura aproximada cuando $p=0{,}1$, $n=100$ y $\alpha=0{,}05$?
\end{ejercicio}
